% ============================================================
% Sección 4
% ============================================================
\section{4 Puntos Críticos y Soluciones de Equilibrio}

\begin{tcolorbox}[enunciado]
\textbf{Clases:}
\begin{itemize}[leftmargin=*]
    \item Sol Equilibrio Curva Integral, Campo Direccional 26 Ago 2026 MCA4
    \item Isoclinas, Análisis Cualitativo 1 MCA4 2Sep26
\end{itemize}
\end{tcolorbox}

\vspace{0.3cm}
\textbf{Responde:}

\begin{enumerate}[label=\arabic*., leftmargin=*]
    \item ¿Qué es una solución de equilibrio?
    
    \vspace{0.3cm}
    \noindent \textbf{Encuentra los puntos críticos o soluciones de equilibrio en cada caso:}
    \vspace{0.2cm}
    
    \item \textbf{EDO, con $y(x)$}
    \[ y^{\prime}=x^{3}-2x^{2}-5x+6 \]

    \item \textbf{EDO (arquetipo) con bifurcación nodo-silla, autónoma (vista en clase)} dada:
    \[ \dot{x}=\mu-x^{2} \]
    \begin{enumerate}[label=(\alph*), leftmargin=*]
        \item $\mu<0$
        \item $\mu=0$
        \item $\mu>0$
    \end{enumerate}

    \item \textbf{Logística}
    \[ \frac{dP}{dt}=rP\left(1-\frac{P}{K}\right) \]
    \begin{enumerate}[label=(\alph*), leftmargin=*]
        \item Soluciones de equilibrio
        \item Qué interpretación físico o biológica tienen las soluciones de equilibrio.
    \end{enumerate}

    \item \textbf{Modelo SIR para Gusanos Informáticos}
    
    Usado en el análisis del gusano Code Red (Zou et al., IEEE 2002):
    
    \textbf{Parámetros:}
    \begin{itemize}[leftmargin=*]
        \item $S(t)$: hosts susceptibles
        \item $I(t)$: hosts infectados
        \item $R(t)$: hosts recuperados
        \item $\beta=\alpha\cdot\eta$: tasa de infección
        \item $\gamma$: tasa de recuperación
    \end{itemize}
    
    \begin{align*}
        \frac{dS(t)}{dt} &= -\beta S(t)I(t) \\
        \frac{dI(t)}{dt} &= \beta S(t)I(t)-\gamma I(t) \\
        \frac{dR(t)}{dt} &= \gamma I(t)
    \end{align*}
    
    \textbf{Conservación:} $S(t)+I(t)+R(t)=N$
    
    \textbf{HINT:} En este caso cada ED tiene que evaluar 0.
\end{enumerate}